\documentclass[a4paper,11pt]{article}

\usepackage{fontspec}
\setmainfont{Carlito}
\usepackage[margin=0.6in, top=0.5in, bottom=0.5in]{geometry}
\usepackage{titlesec}
\usepackage{enumitem}
\usepackage{hyperref}
\usepackage{xcolor}
\usepackage{microtype}
\usepackage{needspace}

\hyphenpenalty=10000 \exhyphenpenalty=10000 \emergencystretch=3em

\definecolor{headernavy}{RGB}{12, 83, 166}
\definecolor{linkblue}{RGB}{0, 102, 204}
\definecolor{muted}{RGB}{96, 102, 110}

\hypersetup{colorlinks=true, linkcolor=linkblue, urlcolor=linkblue}

\titleformat{\section}
  {\color{headernavy}\normalsize\bfseries}
  {}{0em}{\MakeUppercase}[\color{headernavy}\titlerule]

\titlespacing{\section}{0pt}{8pt}{4pt}

\setlist[itemize]{noitemsep, topsep=1pt, leftmargin=1.3em, partopsep=0pt, parsep=1pt}

\pagestyle{empty}
\setlength{\parindent}{0pt}
\linespread{1.02}

% ---- building blocks: bold = organisation/project, italic = role/stack, grey = dates/places ----
% v3.5.3 (10 Nov 2026): employer (and place) on line 1, the title then the dates on line 2, so that `pdftotext -raw` reads every role as
% employer -> title -> dates -> bullets (ats_structure.role_order_problems). The dates were on line 1 before, read before the title.
\newcommand{\role}[4]{% org, location, dates, title
  \textbf{#1}{\color{muted}\,\textperiodcentered\,\textit{#2}}\\
  \textit{#4}\hfill{\color{muted}#3}\par}
\newcommand{\roleNoLoc}[3]{\textbf{#1}\\ \textit{#3}\hfill{\color{muted}#2}\par}
\newcommand{\project}[3]{% name, stack, right-hand tag
  \textbf{#1}\hfill#3\\ {\color{muted}\small\textit{#2}}\par}
\newcommand{\gh}[1]{\href{#1}{\small GitHub}}
\newcommand{\badge}[1]{{\color{headernavy}\small\textbf{#1}}}
\newcommand{\skillgroup}[2]{{\color{headernavy}\textbf{#1}}\\ #2\par\vspace{3pt}}
\newcommand{\gap}{\vspace{5pt}}
\newcommand{\num}[1]{\textbf{#1}}   % result numbers only


\begin{document}

% ================= HEADER =================
\begin{center}
  {\LARGE\bfseries Thiruvel Andagurunathan Pandian} \\[3pt]
  {\color{headernavy}\small\textbf{Data Scientist | Per-gene models for 11,459 genes, validated on held-out data}} \\[3pt]
  {\small
  Bristol, UK\nobreak\ \,\textbar\, \href{mailto:apthiruvel@gmail.com}{apthiruvel@gmail.com}\nobreak\ \,\textbar\, \href{tel:+447741016954}{+44 7741 016954}\nobreak\ \,\textbar\, \href{https://www.linkedin.com/in/thiruvel-a-p/}{https://www.linkedin.com/in/thiruvel-a-p/}\nobreak\ \,\textbar\, \href{https://github.com/Thiruvel-AP}{https://github.com/Thiruvel-AP}\nobreak\ \,\textbar\, \href{https://thiruvel-ap.github.io/PersonalPortfolio/}{https://thiruvel-ap.github.io/PersonalPortfolio/}}\par\vspace{4pt}
  \noindent{} Right to Work in UK \,\textbar\, Available to join Immediately
\end{center}

% ================= PROFILE =================
\needspace{5\baselineskip}\section{Profile}
Data scientist who tests models honestly: trained separate models for 11,459 genes on one GPU, and showed on data they had not seen that they were overfitting. Built an AI assistant on AWS for an MSc dissertation with an AstraZeneca affiliation; it checks its answers against official databases and is tested by about 65 automated checks. Handles the whole data workflow in Python and SQL, from combining raw data to testing and deploying the model (Docker, AWS). Worked freelance building AI and machine-learning tools for a private client (LangChain, Hugging Face). Shipped features in a live finance app for a company\textquotesingle{}s 10,000+ employees, including live stock prices and bank linking.

% ================= SKILLS =================
\needspace{5\baselineskip}\section{Skills}
\skillgroup{Machine Learning, Deep Learning \& Statistics}{Machine learning concepts, statistical analysis (descriptive statistics, correlation analysis), predictive modelling (gradient boosting with XGBoost, random forest, scikit-learn), clustering (K-Means, HAC, DBSCAN, Leiden), deep learning (PyTorch, TensorFlow/Keras, LSTM, BiRNN, CNN), time-series forecasting (cyclic time encoding, TimeSeriesSplit), anomaly detection (Isolation Forest), class-imbalance handling (SMOTE), hyper-parameter tuning (grid search, two-stage Bayesian), model evaluation (held-out, cross-validation, Pearson, MCC), statistical inference and experimental design \textit{(confidence intervals, bootstrap, hypothesis testing; MSc coursework)}, MLflow experiment tracking}
\skillgroup{Python, Data Quality \& Cloud}{Python, pandas, NumPy, SciPy, Matplotlib, Seaborn, Plotly, SQL, PostgreSQL, R (tidymodels), identifier crosswalks, parquet, reproducible pipelines, AWS, Docker, Kubernetes (AWS EKS), serverless deployment (AWS Lambda), containerisation, GitHub Actions CI/CD, FastAPI, TensorFlow Serving, pytest, Git}
\skillgroup{NLP, LLMs, Embeddings \& Evaluation}{NLP (text preprocessing, TF-IDF, SBERT embeddings, text clustering), Hugging Face Transformers, large language models (Claude, OpenAI, open-weight via Ollama), prompt engineering, RAG with embeddings (ChromaDB), agents (LangChain, LangGraph), team LoRA fine-tune of Llama, LLM evaluation (hallucination-rate analysis, deterministic guardrails)}
\skillgroup{Client Delivery \& Communication}{Scoping ambiguous problems with clients, stakeholder demos, working with client engineering teams, code review, transparent reporting of model limitations, root-cause analysis, technical documentation}

% ================= EXPERIENCE =================
\needspace{5\baselineskip}\section{Work Experience}
\needspace{4\baselineskip}
\role{AstraZeneca}{Bristol, UK (Remote)}{May 2026 – Sep 2026}{Data Scientist, MSc Dissertation Project}
\begin{itemize}
  \item Built a generative AI assistant researchers could trust: it checks gene names against official databases and falls back to a local list of 19,213 genes if they go down (LangChain, AWS).
  \item Did the data analysis on 14 research datasets and combined them so the assistant only used trusted data, then put it online on AWS (Python, AWS Lambda).
  \item Found where the assistant went wrong with about 65 automated checks at three levels, fixed failures layer by layer, and reused the checks after changing AI provider (pytest).
  \item Studied graph neural network research and designed how genes and cell lines are linked in one graph, with clustering to find similar ones (Leiden, KNN).
\end{itemize}
\gap
\needspace{4\baselineskip}
\role{180 Degrees Consulting Bristol}{Bristol, UK (Remote)}{Sep 2025 – Dec 2025}{AI Engineer}
\begin{itemize}
  \item Led the design in a five-person team that turned a startup client\textquotesingle{}s vague idea into a working AI prototype in three months; senior leadership rated it 5 stars (Python, MCP).
  \item Designed how a user\textquotesingle{}s question reaches the AI server: one AI model writes the database query and a second runs the statistics before the answer returns (MCP, GraphQL).
\end{itemize}
\gap
\needspace{4\baselineskip}
\role{Freelance}{Chennai, India}{Jan 2025 – Aug 2025}{AI/ML Engineer}
\begin{itemize}
  \item Built a private client an AI finance assistant that reads his income, spending, loans and investments and suggests next steps, self-hosted to keep costs low (Llama, Ollama).
  \item For a private client, built AI chat and document-search (RAG) tools, served machine-learning models as web APIs, and contributed to a team fine-tune of a Llama model (LangChain, FastAPI).
\end{itemize}
\gap
\needspace{4\baselineskip}
\role{Avasoft}{Chennai, India}{Jan 2023 – Sep 2024}{Software Development Engineer Trainee, then Software Development Engineer}
\begin{itemize}
  \item As a trainee, replaced a manual data-collection job with an automatic web scraper, giving the team market trends they used for pricing decisions (JavaScript).
  \item Shipped features in a live finance app as part of a team, connecting live market data and bank accounts through APIs with secure sign-in (Twelve Data, Plaid).
  \item Made the finance app keep working without a signal by saving records on the phone and syncing them in the background (Flutter, Isar).
  \item Built e-commerce and ticketing UIs in Swift/SwiftUI with product search, cart logic, order tracking and booking management.
\end{itemize}

% ================= PROJECTS =================
\needspace{5\baselineskip}\section{Projects}
\needspace{4\baselineskip}
\setbox0=\hbox{\textbf{GTAI Cell Line Finder: multi-omics graph for cancer cell-line selection (PoC)}\hspace{1em}\url{https://github.com/Thiruvel-AP/GTAI-GraphPOC}}%
\setbox2=\hbox{{\small\textit{Python, networkx, Leiden, scikit-learn, RAPIDS cuML}}\hspace{1em}{\small\url{https://github.com/Thiruvel-AP/GTAI-GraphPOC}}}%
\ifdim\wd0>\linewidth
\ifdim\wd2>\linewidth
\textbf{GTAI Cell Line Finder: multi-omics graph for cancer cell-line selection (PoC)}\\ {\color{muted}\small\textit{Python, networkx, Leiden, scikit-learn, RAPIDS cuML}}\\ \hspace*{\fill}{\small\url{https://github.com/Thiruvel-AP/GTAI-GraphPOC}}\par
\else
\textbf{GTAI Cell Line Finder: multi-omics graph for cancer cell-line selection (PoC)}\\ {\color{muted}\small\textit{Python, networkx, Leiden, scikit-learn, RAPIDS cuML}}\hfill{\small\url{https://github.com/Thiruvel-AP/GTAI-GraphPOC}}\par
\fi
\else
\project{GTAI Cell Line Finder: multi-omics graph for cancer cell-line selection (PoC)}{Python, networkx, Leiden, scikit-learn, RAPIDS cuML}{\url{https://github.com/Thiruvel-AP/GTAI-GraphPOC}}
\fi
\begin{itemize}
  \item Trained per-gene regression models for 11,459 genes (OLS and GPU random forest) with resumable chunked training, cutting RAM from 8.2 to 5.5 GB, and validated them on held-out data, surfacing overfitting (train Pearson 0.807 vs test 0.545).
  \item Combined mRNA and protein evidence per gene with Stouffer\textquotesingle{}s method.
\end{itemize}
\gap
\needspace{4\baselineskip}
\setbox0=\hbox{\textbf{Mining Insights from Customer Feedback: NLP \& unsupervised analytics (MSc team project)}\hspace{1em}\url{https://github.com/Thiruvel-AP/TextAnalytics}}%
\setbox2=\hbox{{\small\textit{Python, scikit-learn, SBERT, K-Means, HAC, DBSCAN}}\hspace{1em}{\small\url{https://github.com/Thiruvel-AP/TextAnalytics}}}%
\ifdim\wd0>\linewidth
\ifdim\wd2>\linewidth
\textbf{Mining Insights from Customer Feedback: NLP \& unsupervised analytics (MSc team project)}\\ {\color{muted}\small\textit{Python, scikit-learn, SBERT, K-Means, HAC, DBSCAN}}\\ \hspace*{\fill}{\small\url{https://github.com/Thiruvel-AP/TextAnalytics}}\par
\else
\textbf{Mining Insights from Customer Feedback: NLP \& unsupervised analytics (MSc team project)}\\ {\color{muted}\small\textit{Python, scikit-learn, SBERT, K-Means, HAC, DBSCAN}}\hfill{\small\url{https://github.com/Thiruvel-AP/TextAnalytics}}\par
\fi
\else
\project{Mining Insights from Customer Feedback: NLP \& unsupervised analytics (MSc team project)}{Python, scikit-learn, SBERT, K-Means, HAC, DBSCAN}{\url{https://github.com/Thiruvel-AP/TextAnalytics}}
\fi
\begin{itemize}
  \item Built the text-cleaning and grouping steps for airline support tickets, and traced weak groupings (scores below 0.30) to repeated template wording (Python, scikit-learn).
  \item Showed why airline support tickets would not group by topic (the wording was too similar) but could be sorted well with supervised models; our team project scored 79/100 (Python, scikit-learn).
\end{itemize}
\gap
\needspace{4\baselineskip}
\setbox0=\hbox{\textbf{Banking Customer Churn: propensity model}\hspace{1em}\url{https://github.com/Thiruvel-AP/Banking-Customer-Churn-Prediction}}%
\setbox2=\hbox{{\small\textit{R, tidymodels, XGBoost}}\hspace{1em}{\small\url{https://github.com/Thiruvel-AP/Banking-Customer-Churn-Prediction}}}%
\ifdim\wd0>\linewidth
\ifdim\wd2>\linewidth
\textbf{Banking Customer Churn: propensity model}\\ {\color{muted}\small\textit{R, tidymodels, XGBoost}}\\ \hspace*{\fill}{\small\url{https://github.com/Thiruvel-AP/Banking-Customer-Churn-Prediction}}\par
\else
\textbf{Banking Customer Churn: propensity model}\\ {\color{muted}\small\textit{R, tidymodels, XGBoost}}\hfill{\small\url{https://github.com/Thiruvel-AP/Banking-Customer-Churn-Prediction}}\par
\fi
\else
\project{Banking Customer Churn: propensity model}{R, tidymodels, XGBoost}{\url{https://github.com/Thiruvel-AP/Banking-Customer-Churn-Prediction}}
\fi
\begin{itemize}
  \item Predicted which bank customers were likely to leave, tuning the model in two rounds and scoring it with a measure that is fair to uneven groups (R, XGBoost).
\end{itemize}

% ================= EDUCATION =================
\needspace{5\baselineskip}\section{Education}
\noindent\textit{MSc Data Science}, \textbf{University of Bristol, UK}\unskip\nobreak\hfill\penalty50\hskip1em\null\nobreak\hfill\mbox{Sep 2025 – Sep 2026}\par
\noindent{\small Dissertation with AstraZeneca affiliation (GeneTraceAI)}\par
\noindent\textit{B.E. Electrical \& Electronics Engineering}, \textbf{St. Joseph\textquotesingle{}s College of Engineering}, First Class with Distinction\unskip\nobreak\hfill\penalty50\hskip1em\null\nobreak\hfill\mbox{Sep 2019 – Apr 2023}\par

% ================= CERTIFICATIONS & ACHIEVEMENTS =================
\needspace{5\baselineskip}\section{Certifications \& Achievements}
\begin{itemize}
  \item Model Context Protocol: Advanced Topics, Anthropic, \textit{Jun 2026}
  \item Introduction to Model Context Protocol, Anthropic, \textit{Jun 2026}
  \item AWS Academy Cloud Foundations, \textit{Nov 2025}
  \item 180 Degrees Consulting Data Science Consultant Certificate, \textit{Dec 2025}
  \item Bristol PLUS Award, \textit{Mar 2026}
\end{itemize}

\end{document}